\textbf{Partie A}
\medskip

\begin{enumerate}
\item 1970 - 1960 = 10. Donc on lit $f(10)$.

\begin{center}
\psset{xunit=0.25cm,yunit=5cm,arrowsize=2pt 3,comma,algebraic}
\begin{pspicture*}(-4,-0.1)(33,1.1)
\psaxes[linewidth=1.25pt,Dx=5,Dy=0.5,labelFontSize=\scriptstyle]{->}(0,0)(0,0)(33,1.1)
\psplot[plotpoints=2000,linewidth=1.25pt,linecolor=blue]{0}{33}{1/(1+2.71828^(-0.2*x))}

\uput[dr](0,0.5){A}

\psline[linestyle=dashed,ArrowInside=->](10,0)(10,0.88)
\psline[linestyle=dashed,,ArrowInside=->](10,0.88)(0,0.88)
\uput[d](-1.7,0.95){\red $0,9$}
\uput[d](30,1){\blue $\mathcal{C}_f$}
\end{pspicture*}
\end{center}

Dans la limite de précision du graphique, $f(10) \approx 0,9$.
\item Si $\displaystyle\lim_{t \to + \infty} f(t) = 1$ alors $\mathcal{C}_f$ admet une asymptote horizontale d'équation $y=1$ au voisinage de $+\infty$.
\item $f(0)=\dfrac{1}{2}$ donc $\dfrac{1}{a+1}=\dfrac{1}{2} \iff a + 1 = 2 \iff a = 1$.
\item Soit $m$ ce coefficient directeur. $m =\dfrac{y_B-y_A}{x_B-x_A} = \dfrac{0,5}{10}= 0,05$.
\item 
	\begin{enumerate}
		\item La fonction $f$ est continue et dérivable sur $[0~;~+\infty[$. $f(t)$ est de la forme $\dfrac{1}{u(t)}$, de dérivée $-\dfrac{u'(t)}{u^2(t)}$ Donc en posant $u(t)=a+ \e^{-bt}$, $u'(t)=-b\e^{-bt}$

Donc $\forall t \in [0~;~+\infty[$, $f'(t)=-\dfrac{-b\e^{-bt}}{\left(1+ \e^{-bt}\right)^2}=\dfrac{b\e^{-bt}}{\left(1+ \e^{-bt}\right)^2}$.
		\item Le nombre dérivée de $f$ en 0 est égale à la pente de la droite $(AB)$, tangente à $\mathcal{C}_f$ au point d'abscisse 0. Donc $f'(0)= 0,05$ donc $\dfrac{b}{2^2}=0,05$ donc $b = 4 \times 0,05 = 0,2$.
	\end{enumerate}
\end{enumerate}

\medskip

\textbf{Partie B}

\medskip

On admet, dans la suite de l'exercice, que le taux d'équipement en réfrigérateurs est représenté par la fonction $f$ définie sur l'intervalle $[0~;~+\infty[$ par : $f(t) = \dfrac{1}{1+ \e^{-0,2t}}$.

\begin{enumerate}
	\item $\displaystyle\lim_{t \to + \infty} \e^{-0,2t} = 0$ (par composition). Donc $\displaystyle\lim_{t \to + \infty} f(t)=1$.
	\item La fonction $f$ est continue et dérivable sur $[0~;~+\infty[$, de la forme $\dfrac{1}{u(t)}$, de dérivée $-\dfrac{u'(t)}{u^2(t)}$.

Donc en posant $u(t)=1+ \e^{-0,2t}$, $u'(t)=-0,2\e^{-0,2t}$, $\forall t \in [0~;~+\infty[$, $f'(t)=\dfrac{\e^{-0,2t}}{\left(1+ \e^{-0,2t}\right)^2}$.

$\forall t \in [0~;~+\infty[$, $\e^{-0,2t}>0$ et $\left(1+ \e^{-0,2t}\right)^2>0$ donc $f'(t) > 0$ donc $f$ est croissante sur $[0~;~+\infty[$.
	\item On a $f(0)= \dfrac{1}{2}$ et $\displaystyle\lim_{t \to + \infty} f(t)=1$.

La fonction $f$ est continue strictement croissante sur $[0~;~+\infty[$ à valeurs dans $\left[\dfrac{1}{2}~;~1\right[$. Or $0,97 \in \left[\dfrac{1}{2}~;~1\right[$, donc d'après le corollaire du théorème des valeurs intermédiaires, l'équation $f(x) = 0,97$ admet une unique solution $\alpha$ dans l'intervalle  $[0~;~+\infty[$.
	\item Avec la calculatrice $\alpha \approx 17,4$ donc $\alpha \in ]17~;~18[$.
\end{enumerate}

\medskip

\textbf{Partie C}

\medskip

\begin{enumerate}
	\item $\forall t \in [0~;~+\infty[$, $f(t) = \dfrac{1}{1+ \e^{-0,2t}}$. En multipliant numérateur et dénominateur par $\e^{0,2t}$, on obtient : $f(t)=\dfrac{\e^{0,2t}}{\e^{0,2t}(1+ \e^{-0,2t})}= \dfrac{\e^{0,2t}}{\e^{0,2t} + 1} = \dfrac{\e^{0,2t}}{1 + \e^{0,2t}}$.
	\item $f$ est de la forme $\dfrac{u'(t)}{u(t)}$, de primitive $\ln(u(t))$ (avec $u$ fonction strictement positive car les deux termes du quotient sont strictement positifs) en posant $u(t)=\e^{0,2t}$ et $u'(t)= 0,2\e^{0,2t}$.

En écrivant $f(t)=\dfrac{\e^{0,2t}}{1 + \e^{0,2t}}=\dfrac{1}{0,2}\times \dfrac{0,2\e^{0,2t}}{1 + \e^{0,2t}}=5\times \dfrac{0,2\e^{0,2t}}{1 + \e^{0,2t}}$,

on détermine la fonction $F$, une primitive de $f$ : $F(t)= 5\ln\left(1 + \e^{0,2t}\right)$ pour $t \in [0~;~+\infty[$.
	\item $I = \dfrac{1}{40}\displaystyle\int_0^{40} \dfrac{1}{1 + \e^{0,2t}}\,\text{d}t = \dfrac{1}{40}\displaystyle\int_0^{40} \dfrac{\e^{0,2t}}{1 + \e^{0,2t}} \,\text{d}t =\dfrac{1}{40}\left[5\ln\left(1 + \e^{0,2t}\right)\right]_0^{40}=$

$\begin{gathered}\dfrac{5}{40}\left(\ln\left(\e^8+1\right)-\ln(2)\right)\\=\dfrac{1}{8}\left(\ln\left(\e^8+1\right)-\ln(2)\right)=\dfrac{1}{8} \ln\left(\dfrac{\e^8+1}{2}\right)\approx 0,913\end{gathered}$
\end{enumerate}
